\section{训练数据与后训练}

\subsection{为什么普通语音问答数据不够}

Turn-based SFT\footnote{Supervised Fine-Tuning，监督微调。} 样本通常只有“完整用户音频$\to$完整助手回答”，几乎没有重叠、犹豫、短反馈、抢话、取消和自我修正。模型即使 ASR、SQA\footnote{Spoken Question Answering，口语问答。}、TTS 都很强，也只会等待明确句末。全双工训练样本应是双声道、共享时间轴事件序列
\begin{equation}
D=\{(u_{1:T},a_{1:T},c_{1:T},s_{1:T},q,r)\},
\end{equation}
其中 $c_t$ 是交互动作，$s_t$ 是说话人/重叠状态，$q,r$ 是可选的委派请求与返回。数据不仅要回答“说什么”，还要监督“何时说、何时保持沉默、何时停止”。

真实双声道数据稀缺且隐私敏感。SyncLLM 以 212k 小时合成语音对话和 2k 小时真实对话训练时钟同步能力\cite{veluri2024syncllm}；DuplexChat-Pipe 从公共播客执行语言过滤、下载清洗、diarization、双人片段抽取、speech separation 与 restoration，作者报告产出 282,634 小时英语和 132,723 小时日语数据\cite{nakata2026duplexchat}。这些规模说明数据工程已成为模型能力的主要约束，但自动分离会引入串音、错说话人、伪重叠与恢复伪影，必须保留质量 tag，而不能只报告小时数。

\subsection{建议的渐进训练配方}

公开工作普遍采用渐进式训练。OmniFlatten 依次进行模态对齐、半双工对话学习和全双工对话学习\cite{zhang2024omniflatten}；Freeze-Omni 则冻结文本主干，先接入语音输入输出模态，再进行双工多任务训练\cite{wang2024freezeomni}。一个可复现的六阶段训练方案如下：

\begin{enumerate}
  \item \textbf{Representation pretraining：}优化重建质量、语义保真、流式因果性与低帧率，建立 clean/noisy/echo 条件下的稳定表示。
  \item \textbf{Speech--text alignment：}混合 ASR、TTS、speech continuation、spoken QA，让语音能力对齐文本主干；对文本能力做 replay，避免 catastrophic forgetting。
  \item \textbf{Half-duplex dialogue：}学习指令遵循、角色、语音风格和长上下文，确保“说什么”先收敛。
  \item \textbf{Synthetic duplex conversion：}将文本对话扩展为带时间戳的双声道音频，随机采样停顿、backchannel、重叠、barge-in 与网络延迟。
  \item \textbf{Real duplex adaptation：}用真实双声道对话校准时间分布、韵律和文化差异，降低合成数据的模板偏置。
  \item \textbf{Preference/RL\footnote{Reinforcement Learning，强化学习。}：}针对打断精度、错误抢话、任务完成、说话冗余与恢复能力优化；引入 safety interrupt 和 delegate policy。
\end{enumerate}

损失可写为
\begin{equation}
\mathcal L=\lambda_a\mathcal L_{audio}+\lambda_t\mathcal L_{text}+\lambda_c\mathcal L_{control}
+\lambda_{asr}\mathcal L_{align}+\lambda_s\mathcal L_{safety}+\lambda_d\mathcal L_{delegate}.
\end{equation}
不同 loss 对应的 token 数量级差异很大。audio token 的密度远高于 control token；若简单按 token 平均，模型会优先优化声学细节而忽略稀有但关键的正确打断。因此应对 control/action event 进行重加权，并按场景报告 precision/recall，而非只看总 loss。

\subsection{数据流水线的可量化质量控制}

建议为每个音频片段保留以下质量字段：说话人纯度、重叠比例、信噪比、残余回声、说话人分离置信度、语音分离伪影分数、语言与口音、轮次切换间隔、简短应答标签、隐私与安全标签以及许可来源。关键门槛包括：（1）双声道互相关过高时判定为串音泄漏；（2）ASR 转写与说话人时间轴不一致时降低样本权重；（3）将极端负间隔和长时间重叠单独分桶；（4）在训练集与评测集之间，对同一播客、同一说话人和近重复音频进行去重。

对于合成数据，需要进行\emph{双向校准}：先由真实数据给出轮次切换间隔、重叠时长和简短应答频率的目标分布，合成器再按照这些分布采样；训练完成后，还要检查模型输出是否偏离真实分布。单纯追求更低的响应延迟会诱导模型频繁抢话，因此必须同时把错误打断率作为保护指标。

\section{Stateful Serving：把 audio frame 按时送达}

实时语音 serving 的容量上限不是单位时间完成多少 request，而是多少并发 session 能持续在各自 playout deadline 前生成 audio frame；平均 RTF\footnote{Real-Time Factor，实时因子；通常指处理一段音频所需时间与该段音频时长之比。} 小于 1 仍可能掩盖 burst prefill 造成的连续 deadline miss。长会话还会使 KV cache 随 audio token 持续增长，因此 context compaction 必须移出 live path，并通过 replacement instance 的后台 prefill、并行运行和安全 cut-over 完成无感 handoff。

\subsection{容量单位从 request 转变为 session-frame}

文本 serving 可以用单请求 latency 换 throughput；实时语音却不能让 session 等待长 prefill。

\FOne OpenAI 明确以“多少并发 session 能持续按时处理 frame”衡量容量，而非 GPU request throughput\cite{openai2026continuous}。

设活跃会话集合为 $S$，第 $i$ 个会话 frame 预算为 $\Delta_i$，则更合理的 SLO\footnote{Service-Level Objective，服务等级目标。} 是
\begin{align}
\mathrm{DMR} &= \frac{\sum_{i,t}\mathbf 1[Q_{i,t}+C_{i,t}>\Delta_i]}{\sum_iT_i},\\
\mathrm{GoodSessions} &= \sum_i\mathbf 1[\mathrm{DMR}_i<\epsilon\land P_{99}(J_i)<J_{max}],
\end{align}
这里分别衡量 deadline-miss ratio（DMR）\footnote{Deadline-Miss Ratio，未能在 deadline 前完成的 frame 占比。}和满足 jitter 门槛的并发会话数。即使平均 real-time factor（RTF）低于 1，也不能保证交互体验，因为突发的 context prefill 仍可能导致连续 frame miss。

调度器应采用 deadline-aware continuous batching：按照最近 playout deadline 对 audio frame 排序；限制每批 context prefill 的规模；为 audio decode 预留算力；将后台 context compaction、转写修订和委派模型 prefill 放入低优先级队列。负载过高时，应依次减少非关键遥测、降低 acoustic refinement 强度、延后后台任务，最后才考虑牺牲交互核心的质量。

\subsection{KV cache 与连续上下文}

Transformer-XL 的 segment recurrence 是较早的跨片段记忆方案\cite{dai2019transformerxl}，但实时语音还要求状态与播放时间严格一致。

Transformer 的每层 KV cache 随时间增长。粗略显存为
\begin{equation}
M_{KV}\approx 2LNH_{kv}d_hb,
\end{equation}
其中 $L$ 是层数、$N$ 是累计 token、$H_{kv}$ 是 KV head 数、$d_h$ 是 head dimension、$b$ 是每元素字节。音频高 token rate 与长会话叠加，使 cache 成为 session 容量瓶颈。PagedAttention 通过分页降低动态 cache 的碎片和复制，FlashAttention 通过 IO-aware tiling 降低注意力的 HBM 访问\cite{kwon2023pagedattention,dao2022flashattention}；但它们不自动解决无限上下文。

可采用分层记忆：（1）最近 $W$ 秒保留原始双流 token，保证打断处理和指代消解；（2）将较早的语音压缩为最终转写和事件摘要；（3）为工具结果保留结构化 provenance；（4）把用户偏好写入独立的长期状态。压缩并非无损操作，它会改变 token prefix 并使旧 KV cache 失效。\FOne OpenAI 将 context compaction 移出 live path，在后台完成 replacement instance 的 prefill 后再切换\cite{openai2026continuous}，从而避免主实例因长上下文 prefill 而停顿。

\subsection{无缝实例切换}

有状态 session 无法像普通 HTTP 请求一样任意负载均衡。官方描述的切换过程是：启动并预热 replacement；用压缩后上下文 prefill；新旧实例并行运行；在安全边界 cut over。这个过程同时解决上下文压缩、实例维护和故障恢复。

\begin{figure}[H]
\centering
\begin{tikzpicture}[>=Latex,x=1.35cm,y=0.75cm,font=\small]
\draw[->] (0,0)--(8.6,0) node[right]{时间};
\node[left] at (0,1.8){旧实例}; \node[left] at (0,0.9){新实例};
\draw[very thick,official] (0,1.8)--(6.2,1.8);
\draw[very thick,public] (2.0,0.9)--(8.2,0.9);
\draw[thick,dashed] (2.0,0.65) rectangle (4.8,1.15);
\node[font=\scriptsize] at (3.4,0.9){warm + compacted prefill};
\draw[<->,hypothesis,thick] (4.8,1.35)--(6.2,1.35) node[midway,above,font=\scriptsize]{并行运行};
\draw[dashed] (6.2,0.2)--(6.2,2.2) node[above]{cut over};
\end{tikzpicture}
\caption{有状态推理实例切换。切换点应位于可恢复的声学边界，并对重复/缺失 frame 做去重。}
\end{figure}

工程上还要处理三类竞态条件：工具结果恰好在实例切换期间到达；最终转写对旧文本进行修订；用户在切流边界插话。建议为所有事件附加单调递增的 sequence number 和 context epoch；新实例只接受 epoch 匹配的状态变更，并按照 playout timestamp 对 media frame 去重。

\subsection{异步委派的 serving 设计}

\FOne 为降低复杂任务的首次响应延迟，系统会预先创建并填充 frontier-model session，保持稳定的 session affinity，并利用 prompt caching\cite{openai2026continuous}。委派耗时可分解为
\begin{equation}
T_{delegate}=T_{route}+T_{queue}+T_{prefill\_delta}+T_{reason}+T_{tool}+T_{return}.
\end{equation}
预填充只能减少 $T_{prefill\_delta}$；推理强度、工具模式定义的大小和串行工具调用轮数仍然决定长尾延迟。因此，交互层需要采用渐进式反馈：先用简短应答告知用户任务正在处理，结果到达后再选择合适时机插入；如果用户已经更换话题，则取消后台任务或降低其优先级。

Speculative decoding 可以由小模型先生成 draft，再交给大模型并行验证，在不改变目标分布的前提下减少串行 decoding step\cite{leviathan2023speculative}；但它更适合 cognition path，不能替代 interaction path 中具有 hard deadline 约束的调度。

安全策略也必须跨层一致。GPT-Live 系统卡称安全系统随对话持续检查输入输出，可 steer、interrupt、播放安全响应或结束会话\cite{openai2026safety}。这意味着安全控制不能只位于委派模型末端，否则在长推理返回前，实时层可能已经生成不当语音。
